\section{Kỹ Thuật Bảo Mật Đường Cong Elliptic (Elliptic Curve Security Engineering)}

\subsection{Câu 1. (15 điểm) Lựa Chọn Đường Cong Khi Chịu Áp Lực (Curve Selection Under Pressure)}

\subsubsection{(a) (8 điểm) Tranh Cãi Về "Backdoor" Của Chính Phủ}

Bạn là kiến trúc sư bảo mật cho một ứng dụng nhắn tin mới. Các nhà mật mã học đang tranh luận xem liệu NIST P-256 có chứa ``backdoor'' do chính phủ chèn vào hay không, trong khi Curve25519 cho thấy các thuộc tính an toàn tốt hơn nhưng ít được hỗ trợ phần cứng rộng rãi hơn.

\textbf{i.} Hãy phân tích các mối quan ngại cụ thể đối với các đường cong do NIST quy định:
\begin{itemize}
    \item Một backdoor có thể trông như thế nào?
    \item Kẻ tấn công (hoặc chủ thể chèn backdoor) có thể khai thác backdoor đó mà không cần phá vỡ các bài toán nền tảng (ví dụ: DLP/ECDLP) ra sao?
\end{itemize}

\textbf{ii.} Thiết kế một khung đánh giá rủi ro để chọn giữa P-256 và Curve25519.
\begin{itemize}
    \item Những yếu tố nào nên ảnh hưởng đến quyết định của bạn? (ví dụ: hiệu năng phần cứng, tính tương thích, mức độ tin cậy của nguồn thuật toán, bảo mật dài hạn, v.v.)
\end{itemize}

\textbf{iii.} Đội pháp lý báo rằng một số quốc gia yêu cầu tuân thủ chuẩn NIST để bán cho chính phủ.
\begin{itemize}
    \item Đề xuất một giải pháp vừa xử lý mối lo ngại về backdoor, vừa đảm bảo tuân thủ yêu cầu pháp lý.
\end{itemize}

\begin{center}
    \textbf{Trả lời:}
\end{center}

\textbf{i. Phân tích mối quan ngại về ``Backdoor'' của NIST}

Mối quan ngại về các đường cong do NIST quy định (như P-256) không xuất phát từ một lỗ hổng toán học đã được chứng minh trong chính đường cong, mà xuất phát từ sự thiếu minh bạch trong nguồn gốc của chúng và vụ bê bối liên quan.

\textbf{Các hằng số ``bí ẩn'':} Tài liệu chỉ ra rằng các đường cong tiêu chuẩn của NIST (ví dụ: P-256) có các hằng số (như hằng số $b$ trong phương trình) do NSA cung cấp mà không có lời giải thích rõ ràng về nguồn gốc của chúng. Điều này làm dấy lên nghi ngờ rằng các hằng số này có thể đã được lựa chọn đặc biệt để chứa đựng một điểm yếu tiềm ẩn.

\textbf{Một backdoor có thể trông như thế nào? (Vụ bê bối Dual\_EC\_DRBG):}

Vụ bê bối Dual\_EC\_DRBG (một bộ tạo số ngẫu nhiên dựa trên ECC) là ví dụ điển hình về một backdoor như vậy.

\begin{itemize}
    \item Backdoor này bao gồm một cặp điểm $P$ và $Q$ trên đường cong, trong đó NSA biết một số bí mật $d$ sao cho $Q = d \cdot P$.
    \item Bộ tạo số ngẫu nhiên này được thiết kế để tạo ra các số ngẫu nhiên mà NSA (với việc biết $d$) có thể dự đoán được.
\end{itemize}

\textbf{Cách khai thác (Không phá vỡ ECDLP):}

\begin{itemize}
    \item Kẻ tấn công không cần phá vỡ bài toán ECDLP. Thay vào đó, họ khai thác việc dự đoán được số ngẫu nhiên (nonce).
    \item Trong các thuật toán như ECDSA (thường dùng với P-256), mỗi chữ ký yêu cầu một số ngẫu nhiên bí mật $k$ (nonce).
    \item Nếu kẻ tấn công (NSA) có thể dự đoán được nonce $k$ này (thông qua backdoor trong bộ tạo số ngẫu nhiên như Dual\_EC\_DRBG), họ có thể dễ dàng tính toán ra khóa riêng (private key) $d$ của người dùng. Đây chính xác là lỗ hổng đã xảy ra trong vụ hack Sony PS3, nơi $k$ bị sử dụng lại.
\end{itemize}

\textbf{ii. Khung đánh giá rủi ro (P-256 vs. Curve25519)}

Khi thiết kế khung đánh giá rủi ro cho một ứng dụng nhắn tin mới, chúng ta cần cân nhắc các yếu tố sau:

\begin{center}
\begin{tabular}{|p{3.5cm}|p{4.5cm}|p{4.5cm}|p{3cm}|}
\hline
\textbf{Yếu tố} & \textbf{NIST P-256} & \textbf{Curve25519} & \textbf{Nguồn tài liệu} \\
\hline
Mức độ tin cậy (Nguồn gốc) & Rủi ro cao. Các hằng số ``bí ẩn'' từ NSA. Lòng tin bị xói mòn nghiêm trọng sau vụ Dual\_EC\_DRBG. & Rủi ro thấp. Được thiết kế bởi Daniel J. Bernstein (DJB) với triết lý ``tham số minh bạch'' và ``không có gì giấu giếm''. & \\
\hline
Bảo mật triển khai & Rủi ro cao. Dễ bị tấn công (ví dụ: Invalid Curve Attacks, Small Subgroup Attacks) nếu không kiểm tra (validate) điểm cẩn thận. & Rủi ro thấp. Được thiết kế để ``an toàn hơn khi triển khai'', giúp giảm thiểu rủi ro từ các lỗi triển khai phổ biến. & \\
\hline
Hiệu năng & Hiệu năng tốt, thường được hỗ trợ bởi phần cứng. & Ưu việt hơn. ``Nhanh hơn'' P-256. Lợi ích lớn về tốc độ và ``giảm tiêu thụ pin'', rất quan trọng cho ứng dụng nhắn tin di động. & \\
\hline
Hệ sinh thái (Tính tương thích) & Được chuẩn hóa bởi NIST. Bắt buộc cho tuân thủ chính phủ (theo đề bài). & Rất mạnh. Đã trở thành tiêu chuẩn và được sử dụng rộng rãi trong: Signal, WhatsApp, TLS 1.3, OpenSSH. Đây là yếu tố then chốt cho một ứng dụng nhắn tin. & \\
\hline
\end{tabular}
\end{center}

\textbf{Quyết định:} Đối với một ứng dụng nhắn tin ưu tiên bảo mật và quyền riêng tư, Curve25519 là lựa chọn vượt trội rõ rệt do có các tham số minh bạch, an toàn hơn khi triển khai và hiệu năng cao.

\textbf{iii. Giải pháp cân bằng giữa tuân thủ và bảo mật}

Để vừa tuân thủ yêu cầu pháp lý (sử dụng chuẩn NIST) vừa xử lý mối lo ngại về backdoor (liên quan đến việc dự đoán nonce), giải pháp là sử dụng đường cong P-256 nhưng loại bỏ sự phụ thuộc vào số ngẫu nhiên không đáng tin cậy.

\textbf{1. Tuân thủ đường cong:} Sử dụng đường cong NIST P-256 cho các hoạt động mã hóa (ví dụ: ECDH) và chữ ký (ECDSA) để đáp ứng yêu cầu pháp lý.

\textbf{2. Xử lý mối lo ngại về Backdoor (Nonce):}

\begin{itemize}
    \item Mối lo ngại chính là việc backdoor (như Dual\_EC\_DRBG) cho phép kẻ tấn công dự đoán số ngẫu nhiên $k$ (nonce) dùng trong chữ ký ECDSA.
    \item Giải pháp là sử dụng chữ ký ECDSA với nonce xác định (deterministic nonce) (theo chuẩn RFC 6979).
    \item Cơ chế này vay mượn ý tưởng cốt lõi của EdDSA: thay vì tạo $k$ từ một bộ tạo số ngẫu nhiên (có thể bị chèn backdoor), $k$ (hoặc $r$ trong EdDSA) được tạo ra bằng cách hash (băm) khóa riêng và tin nhắn.
    \item Bằng cách này, chúng ta ``loại bỏ hoàn toàn nhóm lỗ hổng liên quan đến 'randomness'''.
\end{itemize}

\textbf{Kết luận:} Bằng cách sử dụng P-256 (tuân thủ) kết hợp với nonce xác định (an toàn), chúng ta đáp ứng yêu cầu của chính phủ mà không cần phải tin tưởng vào các bộ tạo số ngẫu nhiên do NIST cung cấp.

\subsubsection{(b) (7 điểm) Kịch Bản Tấn Công "Invalid Curve"}

Một nhà nghiên cứu bảo mật phát hiện rằng triển khai ECDH của bạn không kiểm tra điểm đầu vào, làm cho nó dễ bị tấn công invalid-curve.

\textbf{i.} Hãy thiết kế một tấn công cụ thể lợi dụng lỗ hổng này.
\begin{itemize}
    \item Kẻ tấn công có thể trích xuất thông tin gì từ phía nạn nhân (ví dụ: khóa riêng, các bit của khóa, hay phép đo khác)?
\end{itemize}

\textbf{ii.} Phát triển một chiến lược kiểm tra và xác thực input toàn diện để ngăn chặn lớp tấn công này.
\begin{itemize}
    \item Giải pháp của bạn có ảnh hưởng hiệu năng như thế nào?
\end{itemize}

\begin{center}
    \textbf{Trả lời:}
\end{center}

\textbf{i. Thiết kế tấn công ``Invalid Curve''}

Cuộc tấn công này lợi dụng việc nạn nhân (Bob) không kiểm tra xem điểm công khai (public key) nhận được có thực sự nằm trên đường cong tiêu chuẩn đã thống nhất hay không.

\textbf{Các bước tấn công:}

\begin{enumerate}
    \item \textbf{Thiết lập của kẻ tấn công (Alice):} Kẻ tấn công chọn một đường cong hoàn toàn khác (ví dụ: $y^2 = x^3 + ax + b'$). Đường cong này được lựa chọn đặc biệt vì nó là một đường cong yếu mà kẻ tấn công biết cách giải bài toán ECDLP một cách dễ dàng.
    
    \item \textbf{Gửi điểm giả mạo:} Kẻ tấn công gửi một điểm $A$ nằm trên đường cong yếu này cho nạn nhân Bob.
    
    \item \textbf{Tính toán của nạn nhân (Bob):} Bob nhận điểm $A$ và (do thiếu kiểm tra) tiến hành tính toán khóa chung $S = b \cdot A$ (với $b$ là khóa riêng của Bob).
    
    \item \textbf{Lỗ hổng bị khai thác:} Các công thức toán học dùng để ``cộng'' điểm (phép toán cơ bản của $b \cdot A$) thường không sử dụng hằng số $b$ của phương trình đường cong. Do đó, Bob thực hiện phép tính của mình trên đường cong yếu của kẻ tấn công mà không gặp bất kỳ lỗi nào.
\end{enumerate}

\textbf{Thông tin bị trích xuất:}

Khóa chung $S$ mà Bob tính ra sẽ nằm trên đường cong yếu của kẻ tấn công. Vì kẻ tấn công biết cách giải bài toán ECDLP trên đường cong này, chúng không chỉ tìm ra được khóa chung $S$ mà còn có thể suy ra toàn bộ khóa riêng $b$ của nạn nhân.

\textbf{ii. Chiến lược xác thực input (Validation)}

Một chiến lược xác thực toàn diện phải bảo vệ chống lại cả hai cuộc tấn công dựa trên điểm đầu vào được đề cập trong tài liệu: ``Invalid Curve'' và ``Small Subgroup''.

\textbf{Các bước kiểm tra bắt buộc:}

\begin{enumerate}
    \item \textbf{Kiểm tra phương trình đường cong (Chống Invalid Curve):} Khi nhận được một điểm $A(x,y)$, máy chủ phải xác minh rằng các tọa độ này thực sự thỏa mãn phương trình đường cong tiêu chuẩn đã thống nhất: $y^2 = x^3 + ax + b \pmod{p}$. Nếu phương trình không cân bằng, điểm đó phải bị từ chối ngay lập tức.
    
    \item \textbf{Kiểm tra phân nhóm (Chống Small Subgroup):} Sau khi xác nhận điểm nằm trên đường cong, máy chủ phải kiểm tra xem điểm đó có thuộc phân nhóm an toàn (nhóm lớn, có bậc $l$) hay không. Điều này được thực hiện bằng cách tính $l \cdot A$. Nếu kết quả là điểm vô cực ($\mathcal{O}$), thì điểm đó an toàn; nếu không, nó thuộc một phân nhóm nhỏ và phải bị từ chối.
\end{enumerate}

\textbf{Ảnh hưởng hiệu năng:}

\begin{itemize}
    \item Việc thực hiện các bước xác thực này, đặc biệt là bước 2 (Kiểm tra phân nhóm, vốn là một phép nhân vô hướng $l \cdot A$), sẽ làm tăng thêm chi phí tính toán cho mỗi lần trao đổi khóa.
    \item Tài liệu không định lượng chi phí này, nhưng nhấn mạnh rằng các đường cong hiện đại như Curve25519 được ưu tiên sử dụng vì chúng vừa ``nhanh hơn'' vừa ``an toàn hơn khi triển khai''. Thiết kế của Curve25519 vốn đã giảm thiểu các rủi ro tấn công này, làm cho việc xác thực trở nên hiệu quả hơn (hoặc trong một số trường hợp, không cần thiết) so với các đường cong NIST cũ.
\end{itemize}

\subsection{Câu 2. (15 điểm) Phân Tích Lỗ Hổng Triển Khai (Implementation Vulnerability Analysis)}

\subsubsection{(a) (8 điểm) Phòng Thí Nghiệm Tấn Công Kênh Bên (Side-Channel Attack Laboratory)}

Bạn được giao nhiệm vụ kiểm thử triển khai ECDSA trên một thiết bị nhúng để tìm lỗ hổng kênh bên.

\textbf{i.} Hãy thiết kế một tấn công timing chống ECDSA.
\begin{itemize}
    \item Bạn sẽ cố gắng phân biệt giữa hai phép toán nào? (ví dụ: thao tác nhân điểm vs. thao tác bổ sung, hoặc xử lý trường hợp bù trừ khác nhau, v.v.)
\end{itemize}

\textbf{ii.} Phát triển các biện pháp khắc phục vừa giữ được hiệu năng, vừa chống lại tấn công của bạn.
\begin{itemize}
    \item Những kỹ thuật ``constant-time'' cụ thể nào bạn sẽ áp dụng?
\end{itemize}

\begin{center}
    \textbf{Trả lời:}
\end{center}

\textbf{i. Thiết kế tấn công timing chống ECDSA}

\textbf{A. Kiến thức nền về ECDSA}

\textbf{Quy trình ký ECDSA (đơn giản hóa):}
\begin{enumerate}
    \item Chọn nonce ngẫu nhiên: $k$
    \item Tính điểm: $R = k \cdot G$
    \item Lấy tọa độ $x$: $r = R.x \pmod{n}$
    \item Tính: $s = k^{-1} \cdot (H(m) + r \cdot d) \pmod{n}$
    \item Chữ ký: $(r, s)$
\end{enumerate}

\textbf{B. Điểm yếu có thể khai thác}

\textbf{Các phép toán có thời gian khác nhau:}

\textbf{1. Scalar multiplication: $k \cdot G$}
\begin{itemize}
    \item Phụ thuộc vào các bit của $k$
    \item Có thể phân biệt bit 0 vs bit 1
\end{itemize}

\textbf{2. Modular inversion: $k^{-1}$}
\begin{itemize}
    \item Có thể phụ thuộc vào giá trị $k$
    \item Một số thuật toán không constant-time
\end{itemize}

\textbf{3. Conditional branches trong point addition}
\begin{itemize}
    \item Trường hợp đặc biệt: điểm tại vô cực, điểm đối nhau
    \item Tạo ra timing leakage
\end{itemize}


\textbf{C. Thiết kế tấn công timing}

\textbf{Mục tiêu:} Khôi phục nonce $k$ $\rightarrow$ Tính được private key $d$

\textbf{Tấn công 1: Timing leak trong Scalar Multiplication}

\textbf{Bước 1: Hiểu thuật toán scalar multiplication}

Thuật toán ``Double-and-Add'' (không an toàn):

\begin{verbatim}
def scalar_mult_naive(k, G):
    """
    Tinh k*G bang double-and-add
    KHONG AN TOAN - co timing leak
    """
    result = POINT_AT_INFINITY
    temp = G
    
    # Duyet tung bit cua k (tu LSB den MSB)
    for bit in bits(k):
        if bit == 1:           # <- DIEU KIEN RO RANG
            result = result + temp   # Phep cong diem (cham)
        temp = temp + temp     # Phep double (luon thuc hien)
    
    return result
\end{verbatim}

\textbf{Vấn đề:}
\begin{itemize}
    \item Khi bit = 1: thực hiện cả DOUBLE và ADD
    \item Khi bit = 0: chỉ thực hiện DOUBLE
    \item Timing khác nhau $\rightarrow$ Rò rỉ thông tin về các bit của $k$
\end{itemize}
\textbf{Bước 2: Thu thập timing measurements}

\textbf{Setup thí nghiệm:}

\begin{verbatim}
import time
import statistics

def measure_sign_timing(message, device):
    """
    Do thoi gian ky cho mot message
    """
    timings = []
    
    for i in range(1000):  # Thu thap 1000 mau
        start = time.perf_counter()
        signature = device.sign_ecdsa(message)
        end = time.perf_counter()
        
        timings.append(end - start)
    
    return {
        'mean': statistics.mean(timings),
        'stdev': statistics.stdev(timings),
        'samples': timings
    }
\end{verbatim}

\textbf{Bước 3: Phân tích timing để trích xuất bit của $k$}

\textbf{Phương pháp:}

\begin{verbatim}
def analyze_timing_leakage(timings, threshold):
    """
    Phan tich timing de doan cac bit cua k
    
    Gia dinh:
    - Timing cao -> nhieu bit 1 trong k
    - Timing thap -> it bit 1 trong k
    """
    bit_guesses = []
    
    for t in timings:
        if t > threshold:
            bit_guesses.append(1)  # Bit co the la 1
        else:
            bit_guesses.append(0)  # Bit co the la 0
    
    return bit_guesses
\end{verbatim}

\textbf{Bước 4: Khôi phục nonce $k$}

Khi có đủ thông tin về các bit của $k$:

\begin{verbatim}
def recover_nonce_from_bits(bit_pattern, signature, message, pubkey):
    """
    Voi mot so bit cua k da biet, brute-force cac bit con lai
    """
    r, s = signature
    h = hash(message)
    
    # Gia su biet duoc 200 bit tren 256 bit cua k
    known_bits = bit_pattern  # 200 bit da biet
    unknown_bits = 56         # 56 bit chua biet
    
    # Brute-force 2^56 kha nang (kha thi!)
    for candidate in range(2**unknown_bits):
        k_guess = construct_k(known_bits, candidate)
        
        # Kiem tra: r == (k*G).x ?
        R = k_guess * G
        if R.x == r:
            return k_guess
    
    return None
\end{verbatim}

\textbf{Bước 5: Tính private key $d$ từ $k$}

Công thức ECDSA:

$s = k^{-1} \cdot (h + r \cdot d) \pmod{n}$

$\rightarrow d = r^{-1} \cdot (s \cdot k - h) \pmod{n}$

\begin{verbatim}
def recover_private_key(k, signature, message):
    """
    Khi da biet k, tinh private key d
    """
    r, s = signature
    h = hash(message)
    n = CURVE_ORDER
    
    # d = r^(-1)*(s*k - h) (mod n)
    d = (mod_inverse(r, n) * (s * k - h)) % n
    
    return d
\end{verbatim}\textbf{Tấn công 2: Timing leak trong Modular Inversion}

\textbf{Phân biệt phép toán $k^{-1}$}

Thuật toán Extended Euclidean (không constant-time):

\begin{verbatim}
def mod_inverse_unsafe(a, m):
    """
    KHONG AN TOAN - so vong lap phu thuoc vao gia tri a
    """
    old_r, r = a, m
    old_s, s = 1, 0
    
    # So vong lap phu thuoc vao gcd(a, m)
    while r != 0:           # <- So lan lap khong co dinh
        quotient = old_r // r
        old_r, r = r, old_r - quotient * r
        old_s, s = s, old_s - quotient * s
    
    return old_s % m
\end{verbatim}

\textbf{Timing leak:}
\begin{itemize}
    \item $k$ nhỏ $\rightarrow$ ít vòng lặp $\rightarrow$ nhanh
    \item $k$ lớn $\rightarrow$ nhiều vòng lặp $\rightarrow$ chậm
\end{itemize}

\textbf{Khai thác:}

\begin{verbatim}
def attack_modular_inverse_timing(device):
    """
    Do timing de doan kich thuoc tuong doi cua k
    """
    timings = []
    
    for i in range(10000):
        t = measure_sign_timing(f"message_{i}", device)
        timings.append(t)
    
    # Phan loai k nho vs k lon dua tren timing
    small_k_signatures = [sig for sig, t in timings if t < threshold]
    
    # Tan cong nhung chu ky co k nho (de brute-force hon)
    return small_k_signatures
\end{verbatim}

\textbf{Tấn công 3: Branch Prediction Attack}

\textbf{Phân biệt các trường hợp đặc biệt trong point addition}

Point addition với branches:

\begin{verbatim}
def point_add_unsafe(P, Q):
    """
    KHONG AN TOAN - co conditional branches
    """
    # Truong hop dac biet 1: P = O (diem vo cuc)
    if P == POINT_AT_INFINITY:    # <- Branch 1
        return Q
    
    # Truong hop dac biet 2: Q = O
    if Q == POINT_AT_INFINITY:    # <- Branch 2
        return P
    
    # Truong hop dac biet 3: P = Q
    if P == Q:                    # <- Branch 3
        return point_double(P)
    
    # Truong hop dac biet 4: P = -Q
    if P.x == Q.x and P.y == -Q.y:  # <- Branch 4
        return POINT_AT_INFINITY
    
    # Truong hop thong thuong
    return general_point_add(P, Q)
\end{verbatim}

\textbf{Timing leak:}
\begin{itemize}
    \item Các trường hợp đặc biệt có timing khác nhau
    \item Rò rỉ thông tin về cấu trúc của $k$
\end{itemize}
\textbf{D. Kết quả tổng hợp}

\textbf{Kịch bản tấn công hoàn chỉnh:}

\begin{enumerate}
    \item Thu thập 10,000+ chữ ký với timing measurements
    \item Phân tích timing $\rightarrow$ Đoán được $\sim$200/256 bit của $k$
    \item Brute-force 56 bit còn lại ($2^{56} \approx 1$ giờ trên laptop)
    \item Verify $k$ bằng cách kiểm tra $r = (k \cdot G).x$
    \item Tính private key: $d = r^{-1} \cdot (s \cdot k - h)$
    \item \textbf{Game over} - Có thể ký bất kỳ message nào!
\end{enumerate}

\begin{center}
    \textbf{Trả lời:}
\end{center}

\textbf{ii. Biện pháp khắc phục constant-time}

\textbf{A. Nguyên tắc cơ bản}

\textbf{Mục tiêu:} Đảm bảo timing \textbf{không} phụ thuộc vào dữ liệu bí mật

\textbf{3 nguyên tắc vàng:}
\begin{itemize}
    \item \textbf{Không} có conditional branches phụ thuộc vào secret
    \item \textbf{Không} có memory access patterns phụ thuộc vào secret
    \item Mọi đường thực thi phải mất cùng thời gian
\end{itemize}

\textbf{B. Khắc phục 1: Constant-time Scalar Multiplication}

Thuật toán Montgomery Ladder (constant-time):

\begin{verbatim}
def scalar_mult_constant_time(k, G):
    """
    AN TOAN - constant-time scalar multiplication
    """
    R0 = POINT_AT_INFINITY
    R1 = G
    
    # Duyet tung bit cua k (tu MSB den LSB)
    for i in range(255, -1, -1):
        bit = (k >> i) & 1
        
        # LUON thuc hien CA HAI phep toan
        # Khong co branch conditional
        R0, R1 = cswap(bit, R0, R1)  # Conditional swap
        R1 = R0 + R1                  # Add
        R0 = R0 + R0                  # Double
        R0, R1 = cswap(bit, R0, R1)  # Swap back
    
    return R0
\end{verbatim}

Constant-time conditional swap:

\begin{verbatim}
def cswap(condition, a, b):
    """
    Swap a va b neu condition = 1, giu nguyen neu condition = 0
    Khong co branch, chi dung bitwise operations
    """
    # Tao mask: 0 hoac 0xFFFFFFFF
    mask = -condition  # 0 -> 0x00000000, 1 -> 0xFFFFFFFF
    
    # XOR swap voi mask
    temp = mask & (a ^ b)
    a = a ^ temp
    b = b ^ temp
    
    return a, b
\end{verbatim}

\textbf{Lợi ích:}
\begin{itemize}
    \item Timing hoàn toàn độc lập với giá trị $k$
    \item Luôn thực hiện cùng số phép toán
    \item Hiệu năng tốt (overhead $\sim$10-20\%)
\end{itemize}
\textbf{C. Khắc phục 2: Constant-time Modular Inversion}

Sử dụng Fermat's Little Theorem:

\begin{verbatim}
def mod_inverse_constant_time(a, p):
    """
    AN TOAN - constant-time inversion
    Su dung: a^(-1) = a^(p-2) (mod p)
    """
    # Modular exponentiation voi square-and-multiply constant-time
    return mod_exp_constant_time(a, p - 2, p)

def mod_exp_constant_time(base, exp, mod):
    """
    Tinh base^exp (mod mod) theo constant-time
    """
    result = 1
    base = base % mod
    
    # LUON duyet tat ca 256 bit
    for i in range(255, -1, -1):
        result = (result * result) % mod  # Luon square
        
        # Conditional multiply (constant-time)
        bit = (exp >> i) & 1
        mask = -bit
        temp = (result * base) % mod
        result = select(mask, temp, result)  # Chon khong dung branch
    
    return result

def select(mask, a, b):
    """
    Tra ve a neu mask = 0xFFFFFFFF, b neu mask = 0x00000000
    Khong co branch
    """
    return (mask & a) | (~mask & b)
\end{verbatim}

\textbf{Lợi ích:}
\begin{itemize}
    \item Thời gian cố định bất kể giá trị $a$
    \item Không rò rỉ thông tin qua timing
    \item Hiệu năng: $\sim$256 phép squaring + $\sim$128 phép multiply (trung bình)
\end{itemize}
\textbf{D. Khắc phục 3: Constant-time Point Operations}

Point addition không có branches:

\begin{verbatim}
def point_add_constant_time(P, Q):
    """
    AN TOAN - complete point addition formula
    Hoat dong dung cho MOI truong hop ma khong can branch
    """
    # Su dung complete formulas tu Renes-Costello-Batina (2016)
    # Cac cong thuc nay xu ly TAT CA truong hop dac biet
    # ma khong can if-else
    
    X1, Y1, Z1 = P
    X2, Y2, Z2 = Q
    
    # 12 phep nhan field (co dinh)
    t0 = X1 * X2
    t1 = Y1 * Y2
    t2 = Z1 * Z2
    t3 = X1 + Y1
    t4 = X2 + Y2
    t3 = t3 * t4
    t4 = t0 + t1
    # ... (cac buoc con lai)
    
    # Khong co conditional branch nao
    # Luon thuc hien DUNG 12 phep nhan
    
    return (X3, Y3, Z3)
\end{verbatim}

\textbf{Curve25519 - Được thiết kế cho constant-time:}

\textbf{Đặc điểm:}
\begin{itemize}
    \item Complete addition formulas (không cần trường hợp đặc biệt)
    \item Montgomery curve (phép toán x-coordinate only)
    \item Thiết kế sẵn cho constant-time implementation
\end{itemize}

\begin{verbatim}
# Curve25519 x-coordinate only ladder
def curve25519_scalar_mult(k, u):
    """
    Don gian, nhanh, constant-time
    """
    x_1 = u
    x_2 = 1
    z_2 = 0
    x_3 = u
    z_3 = 1
    
    for i in range(254, -1, -1):
        bit = (k >> i) & 1
        x_2, x_3 = cswap(bit, x_2, x_3)
        z_2, z_3 = cswap(bit, z_2, z_3)
        
        # Montgomery differential addition (5M + 4S + 1*a24)
        # Constant-time, khong co branches
        A = x_2 + z_2
        AA = A^2
        # ... (cac buoc con lai)
        
        x_2, x_3 = cswap(bit, x_2, x_3)
        z_2, z_3 = cswap(bit, z_2, z_3)
    
    return x_2 / z_2
\end{verbatim}

\textbf{E. Khắc phục 4: Sử dụng EdDSA thay vì ECDSA}

\textbf{EdDSA deterministic nonce:}

\begin{verbatim}
def ed25519_sign(private_key, message):
    """
    AN TOAN - EdDSA khong can random nonce
    """
    # Hash private key
    h = SHA512(private_key)
    a = h[:32]  # Scalar
    prefix = h[32:]
    
    # Nonce duoc tinh deterministic
    r = SHA512(prefix || message)
    r = r mod L  # L la order cua curve
    
    # Khong co random -> Khong co timing leak lien quan den RNG!
    R = r * G
    
    hram = SHA512(R || public_key || message)
    S = (r + hram * a) mod L
    
    return (R, S)
\end{verbatim}

\textbf{Lợi ích:}
\begin{itemize}
    \item Loại bỏ hoàn toàn lỗ hổng liên quan đến nonce $k$
    \item Không cần RNG an toàn
    \item Deterministic $\rightarrow$ Dễ test
    \item Constant-time implementations có sẵn (libsodium, etc.)
\end{itemize}

\textbf{F. Biện pháp bổ sung}

\textbf{1. Blinding (Làm mù)}

\begin{verbatim}
def scalar_mult_with_blinding(k, G):
    """
    Them random vao k de che giau gia tri thuc
    """
    # Chon random blind r
    r = random_scalar()
    n = CURVE_ORDER
    
    # k' = k + r*n (mod n) = k (mod n)
    k_blinded = (k + r * n) % n
    
    # Tinh voi k', nhung ket qua giong k*G
    result = scalar_mult_constant_time(k_blinded, G)
    
    return result
\end{verbatim}

\textbf{2. Randomized Projective Coordinates}

\begin{verbatim}
def randomize_point(P):
    """
    Bieu dien cung mot diem voi toa do khac nhau
    """
    X, Y, Z = P
    lambda_rand = random_nonzero()
    
    # (X:Y:Z) va (lambda*X:lambda*Y:lambda*Z) bieu dien cung diem
    return (lambda_rand * X, lambda_rand * Y, lambda_rand * Z)
\end{verbatim}

\textbf{So sánh hiệu năng}

\begin{center}
\begin{tabular}{|l|l|l|l|}
\hline
\textbf{Kỹ thuật} & \textbf{Overhead} & \textbf{Độ an toàn} & \textbf{Khuyến nghị} \\
\hline
Naive implementation & 0\% & Không an toàn & Không bao giờ dùng \\
\hline
Montgomery Ladder & +15\% & An toàn & Khuyến nghị \\
\hline
Complete formulas & +20\% & Rất an toàn & Khuyến nghị \\
\hline
EdDSA (Ed25519) & +10\% & Rất an toàn & \textbf{Tốt nhất} \\
\hline
Blinding & +30\% & Defense-in-depth & Bổ sung \\
\hline
\end{tabular}
\end{center}

\textbf{Khuyến nghị cuối cùng}

\textbf{Nên:}
\begin{itemize}
    \item Dùng Ed25519 thay vì ECDSA cho triển khai mới
    \item Sử dụng Montgomery Ladder hoặc complete formulas
    \item Dùng thư viện đã được kiểm định (libsodium, BoringSSL)
    \item Test với tools như dudect để verify constant-time
    \item Kết hợp nhiều lớp bảo vệ (blinding + constant-time)
\end{itemize}

\textbf{Không:}
\begin{itemize}
    \item Tự implement từ đầu
    \item Dùng branches phụ thuộc secret data
    \item Tin rằng ``compiler sẽ tối ưu đúng''
    \item Bỏ qua timing attacks vì ``quá khó khai thác''
\end{itemize}

\textbf{Nguồn tham chiếu:}
\begin{itemize}
    \item Mục 2.3 (Chương 8): ECDSA và EdDSA
    \item Mục 3.2 (Chương 8): Curve25519 an toàn hơn
    \item Ví dụ 2 (Chương 8): PlayStation 3 hack - bài học về ECDSA
\end{itemize}

\subsubsection{(b) (7 điểm) Khủng Hoảng Xác Thực Ed25519 (Ed25519 Validation Crisis)}

Bạn phát hiện hai thư viện Ed25519 phổ biến chấp nhận các chữ kí khác nhau như hợp lệ cho cùng một thông điệp và khóa công khai — dẫn đến mất khả năng tương tác.

\textbf{i.} Khảo sát nguyên nhân của sự không nhất quán này:
\begin{itemize}
    \item Những bước kiểm tra nào ở các triển khai khác nhau có thể không xử lý giống nhau?
\end{itemize}

\textbf{ii.} Phân tích hệ quả an ninh:
\begin{itemize}
    \item Liệu kẻ tấn công có thể lợi dụng những khác biệt đó để tạo ra tấn công thực tế (ví dụ: cho phép ký giả mạo hoặc tận dụng để né kiểm tra xác thực)?
\end{itemize}

\begin{center}
    \textbf{Trả lời:}
\end{center}

\textbf{i. Khảo sát nguyên nhân của sự không nhất quán}

\textbf{Vấn đề:} Hai thư viện Ed25519 khác nhau có thể cho kết quả VALID và INVALID trên cùng một chữ ký.

\textbf{Nguyên nhân:} Một số thư viện ``lỏng lẻo'' (lenient) đã bỏ qua các bước kiểm tra (validation) quan trọng được quy định trong chuẩn RFC 8032 để tối ưu hóa tốc độ.

\textbf{Các bước bị bỏ qua:}
\begin{itemize}
    \item \textbf{Không kiểm tra $S < L$:} Không kiểm tra $S$ có nằm trong phạm vi chuẩn hay không.
    
    \item \textbf{Không kiểm tra điểm trên đường cong:} Cho phép các điểm từ đường cong khác (yếu hơn).
    
    \item \textbf{Không kiểm tra Small Subgroup:} Không kiểm tra các điểm có bậc (order) nhỏ.
    
    \item \textbf{Không kiểm tra Canonical Encoding:} Chấp nhận nhiều cách mã hóa (encoding) cho cùng một điểm.
\end{itemize}

\textbf{ii. Phân tích hệ quả an ninh}

\textbf{1. Signature Malleability (Khi bỏ qua $S < L$):}

Kẻ tấn công có thể thay đổi $S$ thành $S + L$. Chữ ký $(R, S+L)$ mới vẫn được thư viện lỏng lẻo chấp nhận là hợp lệ.

\textbf{Hậu quả:} Hai chữ ký hợp lệ khác nhau cho cùng một giao dịch, dẫn đến hai Transaction ID khác nhau. Trong blockchain, điều này gây ra Consensus Failure (mạng bị chia rẽ).

\textbf{2. Giả mạo chữ ký (Khi bỏ qua kiểm tra Curve/Subgroup):}

\textbf{Invalid Curve Attack:} Kẻ tấn công gửi một điểm trên một đường cong yếu mà chúng biết cách giải. Nếu thư viện bỏ qua kiểm tra, kẻ tấn công có thể giả mạo chữ ký và thậm chí khôi phục khóa riêng của nạn nhân.

\textbf{Small Subgroup Attack:} Cho phép kẻ tấn công thực hiện các cuộc tấn công giới hạn trong các phân nhóm nhỏ.

\textbf{Khuyến nghị khắc phục:}
\begin{itemize}
    \item \textbf{Tuân thủ nghiêm ngặt RFC 8032:} Thực hiện tất cả các bước kiểm tra, không bỏ sót.
    
    \item \textbf{Sử dụng thư viện đã được kiểm định:} Dùng các thư viện uy tín như libsodium.
    
    \item \textbf{Kiểm tra tương thích chéo:} Đảm bảo các thư viện khác nhau trong hệ thống đều cho cùng một kết quả.
\end{itemize}